\paragraph{Wulff 尺度账本的 \texorpdfstring{$2$}{2}-秩障碍与 fold 异常的内禀不可吸收性}

在定理 \ref{thm:conclusion-serrin-wulff-artin-mazur-witt-triple-collapse}
已将归一化 Serrin 几何支塌缩为单点、
而定理 \ref{thm:conclusion-serrin-wulff-scale-ledger-impossibility}
已排除有限生成加法账本之后，
仍可问：即便不要求有限生成性，而仅取有限素数集的局部化，
其 Pontryagin 对偶是否足以承载真实 fold 异常？
下面三个层面——$2$-秩、最小退化层 parity 核、
以及零陪集的 odd-part——同时给出否定回答。

\begin{theorem}[有限素数局部化的 Wulff 账本至多容纳一位内禀奇偶]\label{thm:conclusion-serrin-wulff-localized-parity-rank-one}
设 \(S\subset\PP\) 为有限素数集，记
\[
G_S:=\ZZ[S^{-1}],
\qquad
\widehat G_S:=\Hom_{\mathrm{cont}}(G_S,\TT)
\]
为其 Pontryagin 对偶。则
\[
\dim_{\FF_2}\widehat G_S[2]\le 1,
\qquad
\widehat G_S[2]:=\{x\in \widehat G_S:2x=0\}.
\]
特别地，\(\widehat G_S\) 不含任何子群 \((\ZZ/2\ZZ)^2\)。
因而，任何由有限素数局部化得到的 Wulff 尺度账本，
都至多携带一个内禀二元 parity 自由度。
\end{theorem}
\begin{proof}
离散阿贝尔群 \(G_S=\ZZ[S^{-1}]\) 中，
\(\ZZ\hookrightarrow G_S\) 的余核同构于
\[
G_S/\ZZ
\cong
\bigoplus_{p\in S}\ZZ[1/p]/\ZZ
\cong
\bigoplus_{p\in S}\QQ_p/\ZZ_p.
\]
对 Pontryagin 对偶函子施加左正合性，得到短正合列
\[
0\longrightarrow
\prod_{p\in S}\ZZ_p
\longrightarrow
\widehat G_S
\longrightarrow
\TT
\longrightarrow 0,
\]
其中第一项利用了标准对偶
\(\widehat{\QQ_p/\ZZ_p}\cong \ZZ_p\)
（参见 \cite{MorrisPontryaginDuality1977}）。

核 \(\prod_{p\in S}\ZZ_p\) 作为加法群无非平凡有限挠元：
若 \(x=(x_p)_{p\in S}\) 满足 \(2x=0\)，
则 \(2x_p=0\) 在 \(\ZZ_p\) 中对每个 \(p\)，
而 \(\ZZ_p\) 无挠，故 \(x_p=0\) 对一切 \(p\in S\)。
因此
\[
\left(\prod_{p\in S}\ZZ_p\right)[2]=0.
\]
由短正合列的长正合序列，商映射在 \(2\)-挠子群上的限制为单射：
\[
\widehat G_S[2]\hookrightarrow \TT[2].
\]
而
\[
\TT[2]\cong \ZZ/2\ZZ
\]
仅一维，故
\(\dim_{\FF_2}\widehat G_S[2]\le 1\)。
\end{proof}

\begin{corollary}[window-\texorpdfstring{$6$}{6} 的异常签名不能内化为任何有限素数 Wulff 账本]\label{cor:conclusion-serrin-wulff-window6-anomaly-not-internalizable}
在 \(m=6\) 上，
若保持定理 \ref{thm:xi-window6-anomaly-signature-three-layer-splitting}
的 fiberwise permutation gauge 语义不变，
则异常签名空间的最小完备维数为
\[
(\ZZ/2\ZZ)^3\oplus(\ZZ/2\ZZ)^5\oplus(\ZZ/2\ZZ)^{13}
\cong
(\ZZ/2\ZZ)^{21}.
\]
因此，不存在任何有限素数集 \(S\) 使得该 window-\(6\) 异常签名可忠实内化到 \(\widehat G_S\) 之内。
有限素数局部化的各向异性 Serrin 尺度账本，
与 window-\(6\) 的异常 bit-geometry 在 \(2\)-秩上根本不相容。
\end{corollary}
\begin{proof}
若存在忠实内化，则需有单射
\[
(\ZZ/2\ZZ)^{21}\hookrightarrow \widehat G_S[2].
\]
但由定理 \ref{thm:conclusion-serrin-wulff-localized-parity-rank-one}，
\(\dim_{\FF_2}\widehat G_S[2]\le 1\)，
故不可能容纳 \(21\) 维初等二元子群。
\end{proof}

\begin{theorem}[已审计偶窗的最小退化层自带 Fibonacci 级 parity 核]\label{thm:conclusion-serrin-wulff-minsector-fibonacci-parity-kernel}
对已审计偶窗
\(m\in\{6,8,10,12\}\)，
记
\[
\mathcal S_{m,d_{\min}(m)}
:=\{w\in X_m:d^{\mathrm{bin}}_m(w)=d_{\min}(m)\}
\]
为 bin-fold 的最小退化扇区，并定义其局部 gauge 群
\[
G_{m,\min}
:=\prod_{w\in \mathcal S_{m,d_{\min}(m)}} \mathfrak S_{d_{\min}(m)}.
\]
则有规范同构
\[
G_{m,\min}^{\mathrm{ab}}
\cong
(\ZZ/2\ZZ)^{F_m}.
\]
特别地，
\[
\rank_{\FF_2}G_{m,\min}^{\mathrm{ab}}
=
F_m
\in
\{8,21,55,144\}
\quad
\text{对应 }m=6,8,10,12.
\]
\end{theorem}
\begin{proof}
定理 \ref{thm:gut-bin-min-sector-forces-collision-divergence}
的已审计偶窗前提给出
\[
d_{\min}(m)=F_{m/2},
\qquad
|\mathcal S_{m,d_{\min}(m)}|=F_m.
\]
因 \(m\ge 6\)，故
\(d_{\min}(m)=F_{m/2}\ge F_3=2\)。
而对任意 \(d\ge 2\)，对称群的交换化满足
\[
\mathfrak S_d^{\mathrm{ab}}\cong \ZZ/2\ZZ.
\]
于是
\[
G_{m,\min}^{\mathrm{ab}}
\cong
\prod_{w\in \mathcal S_{m,d_{\min}(m)}} \mathfrak S_{d_{\min}(m)}^{\mathrm{ab}}
\cong
(\ZZ/2\ZZ)^{|\mathcal S_{m,d_{\min}(m)}|}
=
(\ZZ/2\ZZ)^{F_m}.\qedhere
\]
\end{proof}

\begin{corollary}[有限个 Wulff 局部化因子无法吸收已审计偶窗的最小退化层]\label{cor:conclusion-serrin-wulff-finite-localized-product-insufficient}
设
\[
\mathcal L_r
:=\widehat G_{S_1}\times\cdots\times \widehat G_{S_r}
\]
为 \(r\) 个有限素数局部化 Wulff 账本的直积。
若某一审计偶窗 \(m\in\{6,8,10,12\}\) 的最小退化层 parity 核
\(G_{m,\min}^{\mathrm{ab}}\) 可以忠实嵌入 \(\mathcal L_r\)，则必有
\[
r\ge F_m.
\]
亦即：
\(r\ge 8\) (\(m=6\))，
\(r\ge 21\) (\(m=8\))，
\(r\ge 55\) (\(m=10\))，
\(r\ge 144\) (\(m=12\))。
\end{corollary}
\begin{proof}
由定理 \ref{thm:conclusion-serrin-wulff-localized-parity-rank-one}，
每个 \(\widehat G_{S_j}\) 的 \(2\)-秩至多为 \(1\)。故
\[
\dim_{\FF_2}\mathcal L_r[2]\le r.
\]
若
\(G_{m,\min}^{\mathrm{ab}}\hookrightarrow \mathcal L_r\)
忠实成立，则其 \(2\)-秩不能超过 \(\mathcal L_r[2]\) 的 \(2\)-秩。
再由定理 \ref{thm:conclusion-serrin-wulff-minsector-fibonacci-parity-kernel}，
\(\dim_{\FF_2}G_{m,\min}^{\mathrm{ab}}=F_m\)，
从而 \(F_m\le r\)。
\end{proof}

\begin{theorem}[奇因子零陪集障碍]\label{thm:conclusion-serrin-wulff-odd-coset-obstruction}
记 \(\mathfrak W_K\) 为模去平移后的纯尺度 Wulff 射线，
并在分辨率 \(m\) 上按定理 \ref{thm:xi-time-part9x-serrin-ray-fibonacci-semirings}
的半环骨架识别为
\[
\ZZ/(F_{m+2}\ZZ).
\]
记其 multiplicity profile 的非平凡 Fourier 零集为 \(Z_m^{W}\)，并令
\[
A_m^{W}:=Z_m^{W}\cup\{0\}
\]
为相应 annihilator 子群。则
\[
|A_m^{W}|=\gcd(F_{m+2},2^m),
\]
因而 \(|A_m^{W}|\) 必为 \(2\) 的幂。于是：

\emph{若同一 Fibonacci-adic 半环骨架上的某个折叠因子，
其谱零集包含一个加法子群或其陪集，
且该陪集大小的奇部非平凡，
则该因子不可能来自任何纯尺度 Wulff 射线。}
\end{theorem}
\begin{proof}
纯尺度 Wulff 射线在窗口 \(\{0,\dots,2^m-1\}\) 上的 residue-counting
只是区间整数按模 \(F_{m+2}\) 的余类分配（推论 \ref{cor:xi-time-part9x-serrin-ray-exact-capacity}）；
因此在适当循环重标号后，其多重度剖面为
\[
d_m^{W}(r)=q_m+\ind_{I_m}(r),
\qquad
|I_m|=a_m,
\qquad
a_m\equiv 2^m\pmod{F_{m+2}}.
\]
对任意非零频率 \(k\)，常值部分 Fourier 系数消失，故
\[
\widehat{d_m^{W}}(k)
=\sum_{r=0}^{a_m-1}e^{2\pi i kr/F_{m+2}}.
\]
几何级数为零当且仅当
\[
F_{m+2}\mid k a_m,
\qquad
F_{m+2}\nmid k.
\]
因此零频群 \(A_m^W\) 的阶恰为
\[
\gcd(F_{m+2},a_m)=\gcd(F_{m+2},2^m),
\]
后者显然是 \(2\) 的幂。
于是任何带有非 \(2\)-幂零陪集的谱结构，都不可能属于纯尺度 Wulff 射线。
\end{proof}

\begin{corollary}[三轴最小窗 \texorpdfstring{$m=58$}{m=58} 的 odd-part 证书]\label{cor:conclusion-serrin-wulff-m58-odd-part-certificate}
定理 \ref{thm:gut-half-scale-zero-block-covers-stable-space} 给出：当
\(m+2\equiv 0\pmod 6\)
时，零集包含一个大小为 \(F_{(m+2)/2}\) 的大陪集；
对三轴最小窗 \(m=58\)（命题 \ref{prop:crt-235-min-depth}），
\[
F_{30}=832040,
\qquad
|\mathcal Z_{58}|=839415.
\]
由于
\(832040=2^3\times 5\times 20801\)，
其 \(5\)-整除性使奇部非平凡。
于是 \(m=58\) 的大零点陪集在结构上便已排除纯尺度 Wulff 射线解释；
此处不需要比较 \(|\mathcal Z_{58}|\) 与任何其他零集基数，
仅其 odd-part 即已构成否证证书。
\end{corollary}
\begin{proof}
由定理 \ref{thm:gut-half-scale-zero-block-covers-stable-space}，
\[
\mathcal S_{F_{30}}(F_{60})\subset \mathcal Z_{58}
\]
是一个大小为 \(F_{30}\) 的零点陪集。
由定理 \ref{thm:conclusion-serrin-wulff-odd-coset-obstruction}，
纯尺度 Wulff 射线的零 annihilator 阶只能是 \(2\)-幂。
由于
\[
F_{30}=832040\equiv 0\pmod 5,
\]
其 odd-part 非平凡，故不可能来自纯尺度 Wulff 射线。
\end{proof}

\begin{conjecture}[parity--scale 渐近正交原理]\label{conj:conclusion-serrin-wulff-parity-scale-orthogonality}
若定理 \ref{thm:gut-bin-min-sector-forces-collision-divergence}
的已审计结论
\[
d_{\min}(m)=F_{m/2},
\qquad
|\mathcal S_{m,d_{\min}(m)}|=F_m
\]
可从已审计偶窗 \(\{6,8,10,12\}\)
推广到所有充分大的偶数 \(m\)，则定义
\[
r_{\min}(m)
:=
\min\bigl\{r:
\exists\;\text{\(r\) 个有限素数局部化 Wulff 账本之直积，
可忠实承载最小退化层 parity 核}\bigr\},
\]
必满足
\[
r_{\min}(m)\ge F_m,
\qquad
\liminf_{m\to\infty}\frac1m\log r_{\min}(m)\ge \log\varphi.
\]
亦即：
bin-fold 的 anomaly core 与各向异性 Serrin 的有限素数尺度账本之间
存在\emph{指数级不可约的正交性}——
任何试图以固定有限个 Wulff 局部化因子解释全尺度异常的方案，都将失败。
\end{conjecture}
